\section{Limitations and future work}

For fixed positive mixing scale and fixed nonnegative logarithmic action likelihood-ratio bound, \cref{obj:thm:fixed-binary-minimax} establishes minimax mean-squared-error order \(T^{-1}\), where \(T\) is the observed trajectory length, for the binary experiment class in \cref{obj:def:binary-pomdp-class}. The stable-grid estimator supplies a constructive statistical upper bound, while \cref{obj:prop:exhaustive-grid-arithmetic-complexity} characterizes its candidate-list cardinality. These conclusions address distinct questions. A candidate count specifies the size of the search space; runtime additionally depends on the cost of constructing, evaluating, and selecting candidates. Arithmetic-operation counts, bit complexity, and the precision needed to certify objective comparisons therefore require separate analysis. We formulate these requirements as a stable fitting problem.

\begin{definitionv}{def:polynomial-time-handle}[Stable fitting problem \(\mathfrak H_{\mathrm{poly}}\)]
The problem \(\mathfrak H_{\mathrm{poly}}\) is to construct an algorithm with an explicit cost bound in a specified arithmetic or bit-complexity model, including precision requirements, satisfying the following conditions:
\begin{itemize}
\item \textbf{(Input.)} The algorithm starts from the \(3\times3\) Hankel matrix
\[
\bigl(\widehat m_{i+j+1}-\widehat m_{i+j}\bigr)_{i,j=0}^{2},
\]
where \(\widehat m_k\) denotes the common-window empirical weighted reward moment.
\item \textbf{(Rank adaptation.)} It adapts over recurrence ranks zero through three.
\item \textbf{(Stability.)} It enforces transient roots in
\[
\{z:|z|\leq\alpha\},
\]
where \(\alpha\) is the contraction bound.
\item \textbf{(Objective accuracy.)} It returns a stable four-state scalar realization whose seven-moment objective is within \(T^{-1}\) of the global optimum, with a guarantee of this gap, where \(T\) is the observed trajectory length.
\item \textbf{(Uniformity.)} The guarantee is uniform over rank-deficient strata, with positive singular-value margins allowed to vanish.
\end{itemize}
For fixed \(t_0\), the mixing-scale parameter, the exhaustive estimator is specified by a finite candidate list of size \(O(T^{19})\). Its stated computational bound concerns this candidate count; runtime and arithmetic-operation bounds for fitting and selection form part of the computational target.
\end{definitionv}

Here ``starts from'' includes access to the complete input vector \((\widehat m_0,\ldots,\widehat m_6)\), the horizon \(T\), and the contraction bound \(\alpha\); the displayed Hankel matrix is the device used to adapt to the recurrence rank. In the stated finite-grid version, the global optimum is the minimum over the stabilized candidate class \(\mathcal C_M\) of \cref{obj:def:stable-grid-estimator}. A valid output consists of a probability initial vector, a row-stochastic transition matrix with Dobrushin coefficient at most \(\alpha\), and a reward function taking values in \([0,1]\), together with a certificate for the claimed objective gap. Replacing this finite-grid optimum by an optimum over a continuous realization class would be a stronger, separate computational target.

The cost-certified fitting problem \leanref{sym:\mathfrak H_{\mathrm{poly}}}{\ensuremath{\mathfrak H_{\mathrm{poly}}}} in \cref{obj:def:polynomial-time-handle} couples stability with a global objective certificate. Rank adaptation must cover recurrence ranks zero through three, including configurations where positive singular values approach zero. Consequently, a guarantee conditioned on a fixed positive singular-value margin would address a narrower problem. Likewise, a small local optimization residual would require an additional argument connecting it to the stipulated global objective gap. The computational question is whether these certification requirements can be met with explicit cost and precision bounds.

\begin{remarkv}{oeq:polynomial-time-attainment}[Rank-adaptive computational attainment]
A natural computational question is whether \(\mathfrak H_{\mathrm{poly}}\), the cost-bounded stable fitting problem, can implement fitting and lexicographic selection over, or replace, the \(O(T^{19})\)-sized stable-grid candidate list by a rank-adaptive algorithm with explicit runtime, bit-complexity, and precision semantics. The required output is a stable four-state scalar realization whose seven-moment objective is guaranteed to be within \(T^{-1}\) of the global optimum, uniformly over \(\mathcal M_T^{(2)}(t_0,\zeta)\), the class of binary-state and binary-action experiments satisfying the six law conditions. Here \(T\) is the observed trajectory length, \(t_0\) is the mixing-scale parameter, and \(\zeta\) is the logarithmic action likelihood-ratio bound. The question requires uniformity over every rank-deficient Hankel stratum, with positive singular-value margins allowed to vanish.
\end{remarkv}

\Cref{obj:oeq:polynomial-time-attainment} identifies an algorithmic research target rather than an attained complexity guarantee. Its central requirement is global certification across rank-deficient strata: computational savings must preserve the objective guarantee as the effective recurrence rank changes. Resolving this question would supply a cost-controlled implementation of the specified fitting task and clarify the computational requirements for applying stable moment fitting to long trajectories.

Several statistical questions remain within the fixed-binary setting. The seven common-window moments in \cref{obj:def:observable-moments} support the estimator and upper bound in \cref{obj:thm:stable-grid-upper}; determining the minimal sufficient observable window would clarify how much trajectory-local information stationary evaluation requires under the same law conditions. Limit distributions and confidence intervals are further targets. The mean-squared-error guarantees characterize estimation accuracy, while distributional approximations and interval coverage require additional results, particularly for experiments on or near rank-deficient strata.

Application-specific calibration is another direction. The mixing-scale and action-overlap bounds specify the regime of the statistical guarantees. Methods for selecting and substantiating these bounds in a given application, together with sensitivity analyses for their choice, would connect the stated conditions to empirical practice.

Finally, larger or increasing hidden-state cardinalities define a broader statistical extension. Such an analysis would need to characterize how the required observable window and risk bounds depend on cardinality while specifying the corresponding stationarity, overlap, and contraction conditions. The fixed-binary results in \cref{obj:thm:stable-grid-upper,obj:thm:fixed-binary-minimax} provide the benchmark under the six law conditions of \cref{obj:def:binary-pomdp-class}; rates for broader hidden-state classes remain a separate question.
